\section{Introduction}
\label{sec:1}

Time-series analysis (TSA)~\citep{hamilton1994time,box2015time} has long been a central pillar of quantitative finance~\citep{campbell1997econometrics,tsay2010analysis}, providing the mathematical foundation for modeling market dynamics, assessing risk, and pricing assets. Traditionally, TSA in finance was closely associated with forecasting~\citep{hyndman2021forecasting}: predicting the next price movement, the future volatility level, or the impending macroeconomic shift. Modern probabilistic forecasting goes beyond point estimates by quantifying predictive distributions and forecast uncertainty~\citep{gneiting2014probabilistic}. As the financial industry becomes increasingly data-driven and algorithmic, the demand has evolved from point estimation to the comprehensive generation of plausible market scenarios. This paradigm shift has given rise to Financial Time-Series Generation (FTSG), which extends conventional forecasting to include fully synthetic financial sequences~\citep{Time-GAN,wiese2020,naritomi2020data}, missing-data imputation~\citep{yoon2018gain,tashiro2021csdi,wang2024cwgain}, and probabilistic extrapolation of future financial trajectories~\citep{rasul2021timegrad}.

The motivation for FTSG is deeply rooted in the structural limitations of real-world financial data. Unlike text or images, financial markets are characterized by a low signal-to-noise ratio, extreme non-stationarity, and abrupt regime shifts triggered by geopolitical events or macroeconomic policy changes. Furthermore, empirical financial distributions frequently exhibit "stylized facts" such as heavy tails, volatility clustering, and leverage effects~\citep{cont2001empirical}. While historical data contains instances of these phenomena, critical events like market crashes or sudden liquidity droughts are inherently rare. Consequently, discriminative models trained solely on historical records often suffer from overfitting and fail to generalize during crises. FTSG can address this by creating high-fidelity—and, when trained with formal privacy controls, potentially privacy-preserving—synthetic data that can augment training sets, simulate counterfactual stress-testing scenarios, and train robust reinforcement learning agents in simulated environments.

Across the task taxonomy of this survey, we use FTSG to denote the umbrella research domain (Financial Time-Series Generation), while FTSE, FTSI, and FTSS denote the three task-level formulations for extrapolation, imputation, and synthesis, respectively. This terminology is used consistently throughout the paper and is summarized again in the Problem Formulation and Taxonomy and Applications sections.

Generative AI has rapidly entered financial modeling, as illustrated in Figure~\ref{fig:timeline}. The main families include generative adversarial networks (GANs)~\citep{goodfellow2014gan,Time-GAN,wiese2020,naritomi2020data}, denoising diffusion probabilistic models (DDPMs)~\citep{ho2020ddpm,rasul2021timegrad}, Transformers~\citep{vaswani2017attention} and large language models (LLMs)~\citep{brown2020language,jin2024timellm,ansari2024chronos}. These developments, together with their financial adaptations, have produced a fragmented and rapidly expanding body of literature. Zhang et al.~\citep{Survey1} survey LLMs for time series, while Liang et al.~\citep{Survey2} review time-series foundation models. Kim et al.~\citep{Survey3} and Kong et al.~\citep{Survey4} review deep time-series forecasting architectures, and Nie et al.~\citep{FinSurvey1} and Arsenault et al.~\citep{FinSurvey3} review financial-AI applications. These surveys help map parts of the literature, but they leave several task- and evaluation-oriented gaps.

First, the majority of existing surveys remain heavily skewed toward forecasting (extrapolation) and classification, often treating generation merely as an auxiliary data augmentation step rather than a standalone class of problems encompassing synthesis and imputation. Second, the reviews by Zhang et al.~\citep{Survey1} and Liang et al.~\citep{Survey2} are organized primarily around model families and representation-learning components, whereas the reviews by Kim et al.~\citep{Survey3}, Kong et al.~\citep{Survey4}, Nie et al.~\citep{FinSurvey1}, and Arsenault et al.~\citep{FinSurvey3} emphasize architectures or application domains; this organization obscures differences in task formulation, supervision signals, and conditioning contexts. Third, and most critically, existing literature lacks a systematic discussion on evaluation protocols. Assessing a synthetic financial sequence requires moving beyond generic point-wise error metrics (like MSE or MAE) to evaluate cross-sectional dependence, tail-risk preservation, and downstream utility---areas where current surveys offer little standardization or critical synthesis. Finally, there is a scarcity of resource-oriented reviews that comprehensively catalog financial benchmarks, datasets, and the practical biases (e.g., survivorship bias) that plague empirical financial research.

To substantiate these gaps concretely, Table~\ref{tab:related_surveys} positions our survey against the six prior surveys cited above along five explicit dimensions: task coverage (whether each of the three FTSG sub-tasks---extrapolation (FTSE), imputation (FTSI), and synthesis (FTSS)---is treated as a first-class topic), evaluation treatment (whether a multi-dimensional protocol is provided, beyond per-paper metric tables), dataset cataloging (whether FTSG-relevant financial benchmarks are systematically inventoried), number of references in the bibliography (coded approximately from the available versions), and time window of coverage. Our coding suggests that no prior survey jointly covers all three FTSG sub-tasks with a dedicated evaluation protocol and a curated financial dataset catalog; each gap maps to one of the four motivating critiques above and motivates a corresponding contribution of our survey discussed below.

To address these critical gaps, this paper provides a systematic, evaluation-oriented survey of FTSG, focusing on the intersection of generative AI and financial econometrics over the past five years (2021--early 2026). We elevate the discourse from a mere enumeration of literature to a structured problem formulation and benchmarking perspective. Our contributions are fourfold. First, we propose a task-technique two-level taxonomy. By comprehensively reviewing 131 studies published in 2021--early 2026, we develop a hierarchical classification system that decouples the \textit{what} (extrapolation, imputation, and synthesis) from the \textit{how} (statistical, deep discriminative, GAN, diffusion, foundation-model, and controllable multimodal approaches), enabling a clearer mapping of generative techniques to financial use cases. Second, we establish a systematic evaluation framework for FTSG that moves beyond generic reconstruction metrics to encompass distributional fidelity, stylized-facts preservation, downstream utility, and robustness to regime shifts. Third, we compile and categorize public data resources, APIs, and emerging benchmark suites, including FinTSB \citep{hu2025fintsb} and TFB \citep{TFB}, specifically relevant to financial generation, and discuss practical issues in dataset construction. Fourth, we synthesize open challenges in financial validity and identify actionable future directions, including financial-theory-informed generative modeling and causal representation learning.

% 【文生图】 [A flow diagram illustrating the mapping from financial generation tasks (extrapolation, imputation, synthesis) to practical downstream applications (risk warning, data repair, stress testing, simulation). Style: Professional academic flowchart, Information Sciences journal aesthetic, clear directional arrows.]




The remainder of this paper is organized according to the IRDM structure. The Results section presents our core survey findings across six thematic areas: problem formulation and two-level taxonomy, task-specific analyses of extrapolation, imputation, and synthesis, downstream quantitative applications, evaluation protocols for financial realism, and data resources with benchmarks. The Discussion section critically examines pressing open challenges and future research directions, synthesizing the implications for financial generative modeling. Finally, the Methods section details the systematic literature retrieval strategy and PRISMA-compliant screening protocol~\citep{page2021prisma}, inclusion criteria, and coding procedures governing this survey.